\subsection{问题三分析}
问题三关注的是三层食物链在长期尺度上的行为类型，因此模型必须从复杂生态网中抽出一个仍然有代表性的核心结构。`Hastings-Powell` 三层模型并不是要精确复原长江流域的全部营养级，而是把“生产者--中层消费者--顶层捕食者”这一关键链条保留下来，再用承载力 \(K\) 把营养输入、水质改善和边界压力的综合效应压缩为一个可控尺度。这样，系统的长期行为就可以先被归纳为规则有界振荡、明显周期振荡或更复杂的动态区间。

这一问真正要回答的不是某个时刻的数值，而是“系统最后会落到哪一类状态”。当 \(K\) 较低时，链条更容易表现为规则有界振荡，说明资源供给不足以支撑更强的顶层反馈；当 \(K\) 增大时，捕食和增长之间的相互作用会被放大，系统可能转入明显的周期振荡，甚至进一步逼近更复杂的非线性行为。先把这些类型区分清楚，才能给后续问题四提供一个明确的参照：问题三回答的是“有没有复杂化”，问题四回答的是“复杂化是否已经跨入混沌”。

换言之，问题三并不要求把长江流域的真实结构全部装进模型，而是要找一个能够代表系统动力学方向的最小闭环。只要保住顶层捕食、底层供给和中层转移三者之间的非线性关系，就足以判断承载力变化会把系统推向哪类长期行为。问题四沿同一参数轴判别混沌，而不是在不同尺度上各说各话。本问以“行为分类”为主，不必把精力过度放在单个时点的曲线细节上。

同时，\(K\) 的变化并不是为了拟合某条曲线，而是为了识别承载能力变化对长期行为的方向性影响。因此，本问更关注区间上的动力学分型，而不是单点拟合精度；只要能稳定区分“稳定--周期--复杂”三类状态，就已经完成了问题三的任务。

\subsection{问题四分析}
问题四是在问题三的基础上进一步区分“看起来不规则”和“本质上混沌”。仅凭时间序列的起伏幅度，很难判断系统到底是长周期振荡还是混沌，因此需要引入最大李雅普诺夫指数作为判别量。若指数为负，扰动会被压回；若接近零，系统更接近周期轨道；若在扫描网格下持续为正，则说明初值差异会被指数放大，系统进入混沌可观测区间。

这里尤其要避免把某一个临界值写成绝对阈值。由于指数计算依赖步长、初值和网格密度，\(K\approx 3.22\) 更合适的表述应当是“在当前扫描分辨率下，正最大李雅普诺夫指数首次出现的起点区间”。粗扫结果与加密扫描结果共同支持区间性的观测判断，而不是把一个网格点直接升格为硬边界。这样既尊重数值分析的不确定性，也能准确说明混沌是如何在参数轴上逐步显现出来的。

参数扫描分为两步：粗扫定位区间，加密扫确认正指数并缩小范围；如果两者都指向同一段参数，就可以把这段写成“当前分辨率下的混沌起点区间”。这一判定避免把网格点误写成普遍阈值，也提醒读者：这类数值结论依赖采样精度，属于当前模型和当前扫描设置下的观察结果，而不是脱离条件的绝对边界。问题四的核心不是追求一个永远不变的临界数，而是给出足够稳健、可复核的混沌判定口径。

如果需要继续细化，还可以通过更密的网格或不同初值重复验证，以说明结论对数值设置的稳健性；问题四的表述从“某个点的现象”转为“一个可复核的区间判断”。

如果把它放到图上展示，还应在曲线旁标明扫描范围和加密区间，避免读者把单个采样点误读成普适阈值。
